% !TeX root = ../navier-stokes-manuscript.tex
\subsection{The finite piece that reaches a time slice}\label{sub:BodyFiniteAdjacentRectangles}

Let \(I_*=(0,T_*)\) be the interval of the proposed finite maximal history, and
stop first at a strict cutoff \(T<T_*\).  Control of a response one spatial
step above \(H^M\) does not by itself place that response continuously at
\(t=T\).  Its time derivative must approach the same trace from one spatial
step below.

Let \(A=(1-\Delta)^{1/2}\).  For a smooth Hilbert-valued function \(v\),
\begin{equation}\label{eq:BodyAdjacentTraceIdentity}
  \norm{v(b)}_{H^M}^2-
  \norm{v(a)}_{H^M}^2
  =2\operatorname{Re}\int_a^b
  \pair{A^{M-1}v_t(t)}{A^{M+1}v(t)}_{L^2}\,dt.
\end{equation}
Cauchy--Schwarz now places the middle trace between two adjacent controls:
\begin{equation}\label{eq:BodyAdjacentTracePair}
  v\in L^2(0,T;H^{M+1}),
  \qquad
  v_t\in L^2(0,T;H^{M-1}),
\end{equation}
and the conclusion is
\begin{equation}\label{eq:BodyAdjacentTraceConclusion}
  v\in C([0,T];H^M);
\end{equation}
the approximation argument and strong representative are proved in
\Cref{lem:AdjacentHilbertScaleTrace}.

At \(M=2\), the first concrete pair reads
\begin{equation}\label{eq:FirstConcreteAdjacentPair}
  u\in L^2(0,T;H^3),
  \qquad
  u_t\in L^2(0,T;H^1)
  \quad\Longrightarrow\quad
  u\in C([0,T];H^2).
\end{equation}
The velocity lies above the desired trace, its next temporal response lies
below it, and the integral in \eqref{eq:BodyAdjacentTraceIdentity} brings the
two readings to the same time slice.

Apply the same trace to \(V_n\), whose derivative is \(V_{n+1}\).  To carry all
responses through a finite temporal order \(N\) to the common height \(M\),
collect exactly the rows used by those pairings:
\begin{equation}\label{eq:BodyRectangleNorm}
\begin{aligned}
  \mathfrak R_{N,M}(u;T)
  :={}&
  \sum_{n=0}^{N}
  \norm{V_n}_{L^2(0,T;H^{M+1})}^2
  \\
  &+
  \sum_{n=0}^{N}
  \norm{V_{n+1}}_{L^2(0,T;H^{M-1})}^2.
\end{aligned}
\end{equation}
When \(M=0\), the lower edge is read in \(H^{-1}\).  The temporal depth
\(N\) and spatial height \(M\) remain independent.  The number
\(\mathfrak R_{N,M}(u;T)\) measures the size of this finite view of the
history.

The first nontrivial example, \((N,M)=(1,2)\), contains two adjacent traces
that share \(V_1\):
\begin{center}
\captionsetup{hypcap=false}
\captionof{table}{The \((1,2)\) rectangle.  The shared response closes the first trace and begins the second.}
\label{tab:FirstAdjacentRectangle}
\renewcommand{\arraystretch}{1.3}
\begin{tabular}{@{}c c c@{}}
\toprule
response & upper row & lower row \\
\midrule
\(V_0=u\) & \(V_0\in L^2H^3\) & \(V_1\in L^2H^1\) \\
\(V_1=u_t\) & \(V_1\in L^2H^3\) & \(V_2\in L^2H^1\) \\
\bottomrule
\end{tabular}
\end{center}

The scalar norm does not remember which field produced its entries.  Write
\(\cR_{N,M}[u_0;T]\) for this velocity block together with the pressure rows
reconstructed by \eqref{eq:BodyTemporalPressureRecurrence}.  Every entry is a
derivative of the same velocity or of the normalized pressure forced by that
velocity.  The rectangle is a finite view of one fluid history, not a second
object laid over it.

Classical smoothness on a strict subinterval makes each chosen rectangle
finite:
\begin{equation}\label{eq:StrictSubintervalRectangleBound}
  \mathfrak R_{N,M}(u;T)<\infty
  \qquad(T<T_*,\;N,M<\infty).
\end{equation}
The terminal face changes the order of the quantifiers:
\[
  \forall\,T<T_*:\quad \mathfrak R_{N,M}(u;T)<\infty
  \quad\not\Longrightarrow\quad
  \sup_{0<T<T_*}\mathfrak R_{N,M}(u;T)<\infty.
\]
Fixing \((N,M)\) first leaves one demand: the same bound must survive every
cutoff approaching \(T_*\).  At the lowest rectangle, that demand can be
reduced to one exact critical integral.
